% !TeX root = ../navier-stokes-manuscript.tex
% Archived before the 2026-08-17 Section 2 rewrite.
\subsection{The pressure-complete mixed jet}\label{sub:ThePressureCompleteMixedDerivativeTower}

Give the spatial direction its own index.
For \(n\in\Nzero\) and \(\alpha\in\Nzero^3\), define
\begin{equation}\label{eq:BodyMixedJetDefinition}
  J_{n,\alpha}
  :=\partial_t^n\partial_x^\alpha u,
  \qquad
  \Pi_{n,\alpha}
  :=\partial_t^n\partial_x^\alpha p.
\end{equation}
Here \(n\) follows the fixed-coordinate response in time, while \(\alpha\) identifies one spatial derivative inside that response.
Neither index can replace the other.

Apply \(\partial_x^\alpha\) to the temporal recurrence \eqref{eq:BodyTemporalVelocityRecurrence}.
For \(\beta\leq\alpha\), write
\[
  \binom{\alpha}{\beta}
  :=\prod_{j=1}^3\binom{\alpha_j}{\beta_j}.
\]
The temporal and spatial product rules give
\begin{align}
  J_{n+1,\alpha}
  &+
  \sum_{a=0}^{n}
  \sum_{\beta\leq\alpha}
  \binom{n}{a}\binom{\alpha}{\beta}
  (J_{a,\beta}\cdot\nabla)
  J_{n-a,\alpha-\beta}
  +\nabla\Pi_{n,\alpha}
  \notag\\
  &=\nu\Delta J_{n,\alpha},
  \qquad
  \nabla\cdot J_{n,\alpha}=0.
  \label{eq:BodyMixedJetRecurrence}
\end{align}
The pressure coordinate is fixed by the same velocity rows:
\begin{equation}\label{eq:BodyMixedPressureRecurrence}
  -\Delta\Pi_{n,\alpha}
  =\partial_x^\alpha\partial_i\partial_j
  \sum_{a=0}^{n}\binom{n}{a}(V_a)_i(V_{n-a})_j.
\end{equation}
The Riesz normalization in \eqref{eq:PressureNormalization} chooses the pressure row in the decaying whole-space class.
Every term still belongs to the original history.
The binomial coefficients count how the time and space derivatives divide between the transporting and transported factors.
The right side contains \(\Delta J_{n,\alpha}\), and the elliptic equation determines the nonlocal pressure coordinate \(\Pi_{n,\alpha}\).

The complete mixed jet is the family
\[
  \mathcal J[u_0]
  :=\left\{
    J_{n,\alpha},\Pi_{n,\alpha}:
    n\in\Nzero,\ \alpha\in\Nzero^3
  \right\}
\]
together with the recurrences and divergence relations above.
This is one family containing every finite derivative of one pair \((u,p)\).

Now the datum condition from Section~1 enters the Tower.
Equation \eqref{eq:FeffermanDatumSobolevRows} gives
\(V_0(0)=u_0\in H_\sigma^m(\R^3)\) for every finite \(m\).
Suppose \(V_0(0),\ldots,V_n(0)\) belong to every finite Sobolev space.
The pressure recurrence \eqref{eq:BodyTemporalPressureRecurrence}, finite-order
Sobolev multiplication, and the boundedness of the Riesz transforms give
\(Q_n(0)\in H^m\) for every finite \(m\).  Solving
\eqref{eq:BodyTemporalVelocityRecurrence} for \(V_{n+1}(0)\) then puts that
row in every \(H_\sigma^m\) as well.  Induction in \(n\), followed by spatial
differentiation, yields the complete initial jet
\begin{equation}\label{eq:BodyInitialPressureCompleteMixedJet}
  J_{n,\alpha}(0)\in H_\sigma^m(\R^3),
  \qquad
  \Pi_{n,\alpha}(0)\in H^m(\R^3)
\end{equation}
for every finite \(n\), \(\alpha\), and \(m\).  Thus every finite coordinate of
the pressure-complete mixed jet is generated by the single datum in
\eqref{eq:FeffermanDatumClass}.  This initializes the jet at \(t=0\); terminal
control has not yet entered.

The mixed jet has separated the two derivative directions without separating the history.
It locates every finite temporal response and every spatial derivative inside it, but it does not yet show what one more temporal differentiation exposes in space.
That calculation is waiting in the critical-height balance.
